\section{Introducción}

Comprar un alimento empacado suele ofrecer varias presentaciones del mismo tipo. La etiqueta trae números, ingredientes y declaraciones, pero no todos los productos traen los mismos campos. Uno declara azúcar y fibra. Otro no trae el grado de procesamiento. Otro no dice nada sobre alérgenos, es decir, sobre sustancias que una persona necesita evitar. La decisión se toma con ese material, no con una ficha completa.

Comparar esos productos es difícil por tres razones que llegan juntas. Los datos no comparten la misma escala: un nutriente se expresa por 100\,g o 100\,mL, el procesamiento es un grupo, una etiqueta es un sí o un no, y el precio, cuando existe, es un monto. Además, la información no está igual de completa en todos los registros. Y lo que una persona quiere privilegiar puede no coincidir con lo que privilegia otra.

NutriMatch aborda esa comparación para alimentos empacados de México. Toma la información publicada y lo que la persona declara, y devuelve un orden que se puede explicar. Sirve para ver un producto, contrastarlo con otros parecidos y revisar qué dato sostuvo el resultado. La persona decide si se queda con él, si lo cambia o si lo deja fuera.

El sistema no prescribe una dieta, no interpreta síntomas y no observa la compra que ocurre después. Un resultado alto indica mayor compatibilidad con lo que la persona declaró y con la información disponible. No indica que el alimento sea adecuado para cualquier persona.

El documento sigue ese camino. Primero se precisa el problema y qué necesitaría la persona para usar una comparación así. Después se revisa si los datos publicados alcanzaban, qué mostró la exploración y qué decisiones salieron de esos hallazgos. Luego se explica cómo se prepararon los datos, cómo se convierten en un orden y cómo se comprobó que el sistema hace lo que se diseñó para hacer. El cierre vuelve al uso: qué queda resuelto, qué se puede defender y qué no.

Las cifras corresponden al archivo \archivo, del \corte. Si aparece el 19 de septiembre de 2026, es un corte anterior y se señala como histórico.
